% !TeX root = ../main.tex

\chapter{Foundations and Related Work}\label{ch:foundations}

\section{Brain Connectivity and Parcellation Schemes}\label{sec:parcellation}
\outlinenote{Atlases, connectivity estimation, the 200-region choice and what it commits
to.}

\section{Graph Representation Learning in Neuroimaging}\label{sec:graph-learning}
\outlinenote{Graph neural networks on connectomes, graph autoencoders, and pooling
operators.}

\section{Longitudinal and Spatiotemporal Modelling}\label{sec:longitudinal-modelling}
\outlinenote{Recurrent approaches and spatiotemporal graph transformers, including
Brain-TokenGT as the published comparator introduced in \autoref{ch:baselines}.}

\section{Open Problems Addressed in This Thesis}\label{sec:open-problems}
\outlinenote{Four open problems, each mapped to one research question. Closes with a short
paragraph stating the position this thesis takes on each.}

\subsection{Where Spatial Aggregation Belongs in a Longitudinal Pipeline}\label{sec:gap-aggregation}
\outlinenote{Published longitudinal connectome classifiers aggregate a visit's graph to a
single vector before the temporal model sees it, so region identity is discarded before any
trajectory is modelled. Whether that ordering is necessary, and whether the graph topology
carries signal beyond the connectivity matrix it is built from, are open. This is the gap
RQ3 addresses.}

\subsection{What Self-Supervision on Connectomes Learns}\label{sec:gap-ssl}
\outlinenote{Reconstruction is the default pretext task for connectome autoencoders, but
the objective rewards subject idiosyncrasy, and between-patient variation within a cohort
can exceed between-group variation. Whether a reconstruction-pretrained encoder yields
features useful for a downstream group contrast is open. This is the gap RQ2 addresses.}

\subsection{Evidence Standards Under Small Samples}\label{sec:gap-evidence}
\outlinenote{At cohort sizes in the low hundreds, three practices are unsettled: which
floor a model must clear before an architectural claim is admissible, whether same-seed
replicate variance is reported at all, and how longitudinal evaluation is protocolised.
Treated here as one question about what a claim must survive.}

\subsection{Time and Cohort as Confounds}\label{sec:gap-confounds}
\outlinenote{Follow-up duration correlates with the outcome and is rarely controlled;
visit-time semantics differ between cohorts in ways that silently drop or misalign scans;
and pooling cohorts to buy sample size introduces cohort identity as a decodable nuisance
variable. This is the gap RQ4 addresses.}
